\section{Analysis details}\label{sec:report}
\noindent This section summarizes the simulation, reconstruction, optimization, and comparison steps used in this study.

\subsection{Data and Monte Carlo Samples}

\noindent This study used simulated neutron events in the ePIC backward/negative hadronic calorimeter (nHCal). No collision data were used. Geometry descriptions 
were generated from a parameterized DD4hep compact XML template in which the calorimeter was modeled as a 10-layer sampling detector with a 3/3/4 longitudinal segmentation 
pattern. The transverse detector area was fixed to $1\,\mathrm{m}\times1\,\mathrm{m}$ with $10\,\mathrm{cm}\times10\,\mathrm{cm}$ readout cells.

\begin{figure}[htbp]
    \centering
    \includegraphics[width=0.55\linewidth]{plots/detector_segmentation_sketch.png}
    \caption[Detector segmentation sketch]{\it Detector sketch showing the three longitudinal segments, the absorber/scintillator parameterization, and the neutron beam direction used in the study.}
    \label{fig:nhcal_segmentation}
\end{figure}

\noindent The detector response was simulated with DD4hep/DDsim 1.36.0, Geant4 11.4.1, and ROOT 6.38.04\cite{Frank:2018dd4hep,Agostinelli:2003geant4,Allison:2006geant4,Allison:2016geant4,Brun:1997root}. 
The beam configuration used a neutron G4GPS source spanning a continuous momentum range of $0.1$--$2.5~\mathrm{GeV}/c$. The source spectrum was provided as a tabulated user 
histogram whose points approximate an exponential falloff, with linear interpolation between points performed by G4GPS.

\noindent The reference geometry was the uniform nHCal baseline, in which each longitudinal segment uses the same absorber and scintillator thicknesses. Alternative geometries 
were defined by allowing the absorber and scintillator thickness to vary independently in each of the three longitudinal segments while keeping the total detector depth in a 
comparable range. The initial exploration used a 150-point Latin-hypercube sample with 5000 simulated events per geometry. These values were chosen as a practical compromise 
between design-space coverage and simulation cost for the first pass. Selected leading candidates were later rerun with 20000 events to reduce statistical fluctuations in the 
top-ranked region while keeping the higher-stat confirmation stage computationally manageable.

\subsection{Event, Particle and PID Selections}

\noindent Only neutron gun events were used in this study. The generated momentum spectrum was fixed by the G4GPS configuration described above. No track-based reconstruction 
or particle-identification selection was applied, since the study focused on calorimeter hit activity from single-particle simulation.

\noindent At the event-processing stage, an event was counted as valid if the recorded Monte Carlo energy was positive. The main performance metric, detection efficiency, was then defined as
\[
\epsilon_{\mathrm{det}}=\frac{N_{\mathrm{detected}}}{N_{\mathrm{valid}}}.
\]
An event was counted as detected if at least one calorimeter cell passed threshold anywhere in the detector. In addition to detection efficiency, two shower-activity observables 
were recorded: the average number of fired tiles per valid event and the average number of fired layers per valid event.

\subsection{Signal Extraction}

\noindent The hit-level signal extraction was based on energy deposits in the scintillator cells. A cell was counted as fired if its deposited energy exceeded the threshold 
assigned to the corresponding longitudinal segment. A layer was counted as fired if at least one cell in that layer fired. For each processed run, the event loop produced the 
detected-event count, detection efficiency, average fired-tile multiplicity, and average fired-layer multiplicity.

\noindent Thresholds were derived from a MIP-based calibration step. A separate muon calibration script was used to extract a reference MIP scale from the baseline geometry. 
The calibration used 10000 muon events and pooled energy deposits by layer and segment. Segment-specific thresholds for a given geometry were then obtained by linearly scaling the 
reference MIP value and multiplying by the chosen MIP fraction parameter, assuming an approximately linear dependence on scintillator thickness. In the optimization runs reported here, the threshold 
setting used a MIP fraction of 0.5. This MIP-based thresholding approach follows the same general calibration logic used in internal nHCal studies\cite{Kosarzewski:2025nhcalstatus}.

\subsection{Corrections for Detector Effects}

\noindent No detector-level correction procedure was applied in this study. The comparison was entirely geometry-to-geometry at the detector-response level 
under a fixed simulation and thresholding workflow. It is important to note that this study is a relatively idealized simulation study. Because the same beam configuration, 
detector template, threshold definition, and event-processing logic were used for both the baseline and all candidate geometries, the comparison is intended as a controlled 
relative benchmark rather than a fully corrected absolute detector-performance measurement.

\subsection{Surrogate Model and Bayesian Optimization}

\noindent The optimization stage used a surrogate-guided loop. Processed run outputs were converted into a run-level CSV and then compacted to one row per geometry. The compact 
dataset contained the geometry parameters together with the neutron detection efficiency, average fired-tile count, and average fired-layer count. These observables were then used 
to train a regression model that proposed additional candidate geometries for simulation. Candidate ranking used a normalized weighted score built from efficiency, fired-layer 
multiplicity, and fired-tile multiplicity, with the baseline geometry used as the zero-point anchor and the 90th percentile of the initial Latin-hypercube sample used as the reference scale.

\noindent The surrogate model was trained on a per-geometry CSV file. Each row in the CSV corresponds to one geometry and contains the 
geometry parameters together with averaged performance metrics. The surrogate input features were:
\begin{itemize}[nosep]
    \item absorber thicknesses in segments 1, 2, and 3
    \item scintillator thicknesses in segments 1, 2, and 3
\end{itemize}
The target variables were:
\begin{itemize}[nosep]
    \item neutron detection efficiency
    \item average fired-tile count
    \item average fired-layer count
\end{itemize}

\noindent The regression model was implemented with LightGBM\cite{Ke:2017lightgbm}. Because multiple target observables were predicted simultaneously, the code wrapped the base LightGBM regressor in a 
multi-output regression interface. The current configuration used 500 estimators, a learning rate of 0.05, 31 leaves, and 0.8 subsampling in both events and features. By default, 
20\% of the compact dataset was held out for validation, and mean absolute percentage error was printed separately for each predicted target, consistently within 2\% for each target. These 
hyperparameters were fixed settings and were not separately tuned in this study.

\noindent Candidate generation for later iterations was handled by quasi-random sampling in the allowed design space. The active workflow sampled a pool of 20000 candidates per 
iteration and generated 5 proposed geometries per Bayesian-optimization (BO) batch. Design constraints were applied first in geometry space, after which the sampled candidates were converted into surrogate 
feature rows and scored with the trained model. The pool size and the batch size were chosen as a practical balance between broad surrogate-guided exploration and downstream simulation 
throughput.

\noindent The surrogate ranking used the normalized weighted score implemented in the scoring helper. The three predicted metrics were neutron detection efficiency, average 
fired-layer count, and average fired-tile count, with weights of 0.5, 0.3, and 0.2, respectively. For each metric, the normalization used the baseline geometry as the zero-point 
anchor and the 90th percentile of the initial Latin-hypercube sample as the reference scale. Values below baseline received no credit. Values between the baseline and the 
90th percentile were scaled linearly. These weights intentionally prioritized neutron detection efficiency as the primary objective, while the larger secondary weight on fired-layer multiplicity 
was meant to give extra preference to longitudinal shower development in a study focused on longitudinal segmentation. The smaller fired-tile weight was retained as a weaker measure of overall shower activity. The 90th-percentile anchor was used to reduce sensitivity to the noisiest upper tail of the initial scan. Values above the 90th percentile were allowed to continue increasing 
with a reduced slope of 0.2, so that the top end of the score did not saturate immediately while still remaining less sensitive to extreme fluctuations. Candidate selection within a BO batch 
also enforced a minimum normalized geometry-space separation of 0.05 between chosen proposals, again as a practical diversity setting rather than a separately optimized parameter.

\noindent The overall workflow was iterative. After a sweep was simulated and processed, the processed outputs were compacted to one row per geometry, and used to retrain 
the surrogate model by appending it onto the CSV from the previous iteration. The updated model was then used to generate the next BO sweep. In parallel with this loop, 
the compact CSV was also rescored directly to identify the current best observed geometry under the same objective used for candidate ranking.

\subsection{Systematic Uncertainties}

\noindent This note did not perform a full detector-systematics study. The dominant uncertainty considered here was Monte Carlo statistical variation between geometry evaluations. 
This was especially important because the efficiency landscape across the scanned geometries was shallow. In the initial 150-point scan, neutron detection efficiencies covered only 
a narrow range, approximately 0.51 to 0.55. This overall efficiency scale is broadly consistent with internal nHCal performance estimates\cite{BrandenburgRiedl:2025nhcalpdr}. Under those conditions, low-statistics evaluations can mis-rank near-top candidates even when the broader landscape is still identified correctly.

\noindent To address this, the study used a staged event-count strategy. The initial Latin-hypercube scan used 5000 events per geometry to map the design space at manageable cost. 
Top-performing candidates were then rerun with 20000 events to test whether their apparent advantage persisted at higher statistics. Several geometries that looked competitive in 
the lower-stat scan did not remain leaders after these higher-stat confirmation runs. This behavior should be treated as a limitation of the initial low-stat ranking for near-tied 
candidates, but not as evidence that the initial scan failed to identify the relevant performance region.

\noindent For the final baseline-versus-optimum comparison, both the baseline geometry and the optimum were evaluated with ten 5000-event runs each. This provides matched repeated-run 
estimates for the main observables and their standard errors on the mean. A broader treatment of systematic variations, such as alternative threshold settings, physics lists, 
or beam-shape variations, is outside the scope of the present note. The number of repeated runs was chosen as a computationally affordable way to estimate run-to-run variation in the final matched comparison.

\subsection{Results}

\noindent The optimization landscape was found to be shallow. The best-performing geometry was already present in the initial 150-point Latin-hypercube sample. 
Later Bayesian-optimization proposals produced candidates with comparable performance, but they did not consistently outperform this geometry once higher-stat reruns were 
used to check the leaders.

\begin{figure}[htbp]
    \centering
    \begin{subfigure}[t]{0.48\linewidth}
        \centering
        \includegraphics[width=\linewidth]{plots/efficiency_hist.png}
        \caption[Efficiency landscape]{\it Histogram of neutron detection efficiencies across the initial 150-geometry scan.}
        \label{fig:efficiency_histogram}
    \end{subfigure}
    \hfill
    \begin{subfigure}[t]{0.48\linewidth}
        \centering
        \includegraphics[width=\linewidth]{plots/sorted_efficiency.png}
        \caption[Sorted efficiency scan]{\it Sorted neutron detection efficiencies from the initial 150-geometry scan, with the baseline shown for comparison.}
        \label{fig:sorted_efficiency_scan}
    \end{subfigure}
\end{figure}

\noindent Figures~\ref{fig:efficiency_histogram} and \ref{fig:sorted_efficiency_scan} provide complementary views of the same shallow landscape. The histogram shows that the scanned geometries occupy only a narrow efficiency range, while the ranked plot shows that many geometries remain close in performance and that the baseline lies within the broader spread of the sampled designs.

\noindent The final matched comparison used the uniform baseline geometry and the optimum, each evaluated in ten independent 5000-event runs. The quoted uncertainties in this 
paragraph are standard errors on the mean. The baseline achieved a neutron detection efficiency of $0.53172 \pm 0.00331$, while the optimum achieved $0.54184 \pm 0.00276$, 
corresponding to an absolute gain of 0.01012 and a relative gain of about 2\%. Over the same comparison, the average fired-tile multiplicity increased from $1.09068 \pm 0.00897$ 
to $1.13452 \pm 0.00660$, and the average fired-layer multiplicity increased from $0.97710 \pm 0.00786$ to $1.01348 \pm 0.00553$. These correspond to increases of about 4.4\% and 4.1\%, 
respectively.

\noindent Tables~\ref{tab:baseline_geometry}--\ref{tab:rerun_b_geometry} summarize the baseline, the optimum, and the next two best-performing geometries among the higher-stat reruns. All thickness values are given in cm.

\begin{center}
    \begin{minipage}[t]{0.47\textwidth}
        \centering
        \begin{tabular}{lc}
            \toprule
            Metric & Value \\
            \midrule
            $t_{\mathrm{abs},1}$ & 4.0000 \\
            $t_{\mathrm{scin},1}$ & 0.4000 \\
            $t_{\mathrm{abs},2}$ & 4.0000 \\
            $t_{\mathrm{scin},2}$ & 0.4000 \\
            $t_{\mathrm{abs},3}$ & 4.0000 \\
            $t_{\mathrm{scin},3}$ & 0.4000 \\
            Efficiency & $0.53172 \pm 0.00331$ \\
            Tiles Mean & $1.09068 \pm 0.00897$ \\
            Layers Mean & $0.97710 \pm 0.00786$ \\
            \bottomrule
        \end{tabular}
        \captionof{table}[Baseline geometry]{\it Baseline geometry.}
        \label{tab:baseline_geometry}
    \end{minipage}
    \hfill
    \begin{minipage}[t]{0.47\textwidth}
        \centering
        \begin{tabular}{lc}
            \toprule
            Metric & Value \\
            \midrule
            $t_{\mathrm{abs},1}$ & 3.6138 \\
            $t_{\mathrm{scin},1}$ & 0.5121 \\
            $t_{\mathrm{abs},2}$ & 3.6732 \\
            $t_{\mathrm{scin},2}$ & 0.4663 \\
            $t_{\mathrm{abs},3}$ & 4.0268 \\
            $t_{\mathrm{scin},3}$ & 0.4801 \\
            Efficiency & $0.54184 \pm 0.00276$ \\
            Tiles Mean & $1.13452 \pm 0.00660$ \\
            Layers Mean & $1.01348 \pm 0.00553$ \\
            \bottomrule
        \end{tabular}
        \captionof{table}[Optimum geometry]{\it Optimum geometry.}
        \label{tab:optimum_geometry}
    \end{minipage}

    \vspace{0.75cm}

    \begin{minipage}[t]{0.47\textwidth}
        \centering
        \begin{tabular}{lc}
            \toprule
            Metric & Value \\
            \midrule
            $t_{\mathrm{abs},1}$ & 3.6043 \\
            $t_{\mathrm{scin},1}$ & 0.5100 \\
            $t_{\mathrm{abs},2}$ & 3.6071 \\
            $t_{\mathrm{scin},2}$ & 0.3395 \\
            $t_{\mathrm{abs},3}$ & 3.5849 \\
            $t_{\mathrm{scin},3}$ & 0.3072 \\
            Efficiency & 0.53705 \\
            Tiles Mean & 1.14130 \\
            Layers Mean & 1.02010 \\
            \bottomrule
        \end{tabular}
        \captionof{table}[High-stat candidate A]{\it First higher-stat comparison candidate.}
        \label{tab:rerun_a_geometry}
    \end{minipage}
    \hfill
    \begin{minipage}[t]{0.47\textwidth}
        \centering
        \begin{tabular}{lc}
            \toprule
            Metric & Value \\
            \midrule
            $t_{\mathrm{abs},1}$ & 3.5150 \\
            $t_{\mathrm{scin},1}$ & 0.5771 \\
            $t_{\mathrm{abs},2}$ & 3.5268 \\
            $t_{\mathrm{scin},2}$ & 0.4571 \\
            $t_{\mathrm{abs},3}$ & 4.0344 \\
            $t_{\mathrm{scin},3}$ & 0.4120 \\
            Efficiency & 0.53670 \\
            Tiles Mean & 1.14450 \\
            Layers Mean & 1.02065 \\
            \bottomrule
        \end{tabular}
        \captionof{table}[High-stat candidate B]{\it Second higher-stat comparison candidate.}
        \label{tab:rerun_b_geometry}
    \end{minipage}
\end{center}


\subsection{Summary and Conclusion}

\noindent This study optimized the absorber and scintillator thicknesses in a simplified 10-layer, three-segment model of the ePIC nHCal using simulated neutron events. The optimization used a surrogate-guided workflow trained on DD4hep/DDsim and Geant4 detector-response outputs, with candidate geometries ranked by a weighted combination of neutron detection efficiency, fired-layer multiplicity, and fired-tile multiplicity.

\noindent The best geometry improved the neutron detection efficiency from $0.53172 \pm 0.00331$ for the uniform baseline to $0.54184 \pm 0.00276$, corresponding to a relative gain of about 2\%. The same comparison showed increases of about 4\% in both fired-tile and fired-layer multiplicities. These results suggest that nonuniform longitudinal segmentation can provide a modest improvement in low-energy neutron response.

\noindent The observed efficiency landscape was shallow, and several near-leading candidates changed rank after higher-statistics reruns. The result should therefore be interpreted as a controlled simulation indication rather than a final detector-design recommendation. Further confirmation with larger event samples, threshold variations, and more realistic detector and beam conditions would be needed before drawing stronger design conclusions.

\subsection{Discussion}

\noindent The efficiency gain is modest. However, the improvement is supported by matched repeated-run statistics, and the multiplicity observables show a clearer separation between the baseline and the optimum. The study therefore suggests that nonuniform longitudinal redistribution can improve low-energy neutron response in the ePIC nHCal, but only within a fairly shallow performance landscape. In this setting, the surrogate-guided workflow was useful for reducing the search burden and identifying promising regions, but higher-stat confirmation remained necessary to distinguish the best few candidates reliably.
